\documentclass{otengineering}

\projectname{Modular Magnetic Headphones}
\projectid{ENG-001}
\projectversion{0.1}
\projectstatus{\statusconcept}
\projectcategory{Consumer Electronics / Hardware Design}
\projectstarted{17 May 2025}
\projectupdated{\today}

\begin{document}

\makeotengtitle%

\section{Project Overview}

\begin{idea}
    Design a semi-modular headphone system with retractable magnetic headbands,
    power/data-sharing connectors, ANC synchronisation, and an interactive earcup screen.
\end{idea}

\section{Engineering Questions}

\begin{questionbox}
    Can magnetic connectors safely support mechanical attachment, charging, and data transfer
    without compromising comfort or reliability?
\end{questionbox}

\section{Design Decisions}

\begin{decision}
    \entrydate{\today}
    \textbf{Decision:} Explore MagSafe-style connector points between the retractable band and earcups.

    \reason{This may simplify attachment while supporting modularity.}

    \tradeoff{Magnetic alignment may introduce cost and durability challenges.}
\end{decision}

\section{Experiments}

\begin{experiment}
    \entrydate{\today}
    \textbf{Objective:} Test whether a retractable rail can fit inside the earcup housing.

    \textbf{Method:} Begin with hand sketches, then produce a simplified 3D model.

    \successrating{4}
    \confidencerating{5}
    \pursuerating{Yes}
\end{experiment}

\section{Failures and Problems}

\begin{failure}
    \entrydate{\today}
    \textbf{Problem:} The first headband concept may require more internal space than expected.

    \observation{The earcup may become too bulky if the rail mechanism is not compact.}

    \nextstep{Compare telescopic, flexible-band, and folding-arm mechanisms.}
\end{failure}

\section{Discoveries}

\begin{discovery}
    \entrydate{\today}
    \context{Thinking about modular headphones.}

    \observation{A magnetic connector might serve three roles: alignment, attachment, and electrical contact.}

    \textbf{Potential Applications:}
    \begin{itemize}
        \item Modular headbands
        \item Swappable earcups
        \item Charging docks
    \end{itemize}
\end{discovery}

\section{Sketches}

\sketchbox{Draw the retractable headband concept here.}

\section{Engineering Calculation}

\begin{calculation}
    \calcfield{Purpose}{Estimate the holding force needed for the magnetic connector.}
    \calcfield{Inputs}{Estimated headphone mass, expected movement, contact area.}
    \calcfield{Assumptions}{The user will experience normal walking and head movement only.}
    \calcfield{Calculation}{To be added.}
    \calcfield{Interpretation}{To be added after first estimate.}
\end{calculation}

\section{Future Research Queue}

\begin{researchqueue}
    \textbf{Investigate:}
    \begin{itemize}
        \item Hall-effect sensors
        \item Magnetic pogo pins
        \item Flexible PCB connectors
        \item Bluetooth antenna placement
    \end{itemize}
\end{researchqueue}

\section{Note to Future Self}

\begin{futurememo}
    Do not redesign the entire headphone too early. First prove the retractable mechanism,
    then test power/data connection, then think about aesthetics.
\end{futurememo}

\section{What I Wish I Knew Earlier}

\begin{wisdom}
    Mechanical feasibility should be checked before adding complex electronics.
    A beautiful product idea can fail if the housing cannot physically contain the mechanism.
\end{wisdom}

\end{document}